% ----------------------
\subsection{Section 19 (Circa Summer 1829)}
% ----------------------
	\label{DC19}
		
	% -----
	\subsubsection{Note on the JSP and BC versions}
	% -----
	
		The version featured in the JSP comes from BC, so for this section they are the same; the BC is therefore not included in the alternate versions below. Some of this section does appear in RB1; there is a spot for it below each verse, and you can read the entire thing \hyperref[EXSS-DC19-RB1]{here}. The RB1 account picks up in the middle of \hyperref[DC19-20]{verse 20} below, and continues until the end of the section.
	
	% -----
	\subsubsection{Historical Background}
	% -----
		\label{DC19-HB}
		
		% --
		\paragraph{JSP}
		% --
		
			``This revelation, directed to Martin Harris, clarified doctrines regarding the nature of God, Christ's atonement, repentance, and the afterlife, and it counseled Harris on a variety of matters. Its immediate purpose was to assure payment to printer E. B. Grandin by commanding Harris to `not covet thine own property, but impart it freely to the printing of the book of Mormon.'
			
			``Although the first two published versions date this revelation to March 1830, the earliest extant manuscript and the context for the revelation better support a summer 1829 date.%
				\footnote{%
					% \label{}%
					Another point to make in favor of 1829 involves the discussion of ``endless torment'' near the beginning of this section. Besides its one occurrence here, that phrase appears seven times in the scriptures, and all of them are in the BOM. It could be that JS was wondering about that idea and the Lord decided to explain it, or to give some light on it at least, in this revelation.
					}
			The earliest manuscript version, found in Revelation Book 1, is only partially extant. It includes only the final portion of the revelation and does not bear a date. But when John Whitmer inscribed it in Revelation Book 1, he placed it within a series of 1829 revelations and listed it in the book's index as an 1829 revelation. When the revelation was first printed, in the 1833 Book of Commandments, it was dated March 1830, with the location and date in parentheses. The Book of Commandments used parentheses frequently to supply information not included in manuscript sources, while the 1835 edition of the Doctrine and Covenants eliminated all such parentheses, with one exception: it retained the parentheses around the dating of this revelation, suggesting that the 1830 date was not certain. Although a recollection by Joseph Knight Sr. seems to provide evidence for a March 1830 date, Knight's placement of the revelation in the later time frame is likely explained by his reliance on the 1835 Doctrine and Covenants.%
				\footnote{%
					% \label{}%
					JSP note: ``In March 1830, Knight traveled with JS to Manchester, New York. As they arrived, they met Martin Harris, who was distraught because no one wanted to buy the Book of Mormon. According to Knight's later narrative, Harris told JS, `I want a Commandment why says Joseph fullfill what you have got But says he I must have a Commandment.' That night Harris and Knight stayed at the Smith home, and when Harris departed the next morning, Knight heard him again tell JS that `he must have a Commandment.' When reconstructing this episode in his narrative, Knight consulted the 1835 Doctrine and Covenants to find the revelation that Harris had demanded and found the printed text of this revelation with the expected date. He then wrote, `And along in the after part of the Day Joseph and Oliver [Cowdery] Received a Commandmant whi[c]h is in Book of Covenants Page 174,' thus associating this revelation, dated `(March 1830)' in the Doctrine and Covenants, with his remembered experience. This passage from Knight's narrative may be viewed as corroborating the March 1830 date. More likely, however, Knight's recollection of the conversation was accurate but he was mistaken in assuming that JS received a new revelation for Harris. By this line of reasoning, when JS told Harris in March 1830 to `fullfill what you have got,' he was referring to the revelation featured here, which Harris had received the previous summer. The March 1830 date requires interpreting the revelation as chastising Harris for delay in selling off enough of his mortgaged property to come up with some or all of the \$3,000 owed to Grandin by the terms of the 25 August 1829 agreement. Given the actual terms of the agreement and the use Grandin made of it, this scenario seems unlikely. This revelation more closely fits a summer 1829 context, and it likely motivated Harris to complete the 25 August 1829 agreement with Grandin soon after.''
					}
			
			``In June 1829, before this revelation was dictated, Harris and JS talked with several printers in Palmyra and Rochester, New York, about printing the Book of Mormon, finally settling on E. B. Grandin of Palmyra. John H. Gilbert, the compositor who assisted Grandin in estimating the cost of the project and later typeset the Book of Mormon, recalled that Harris initiated the negotiations and planned to pay for the printing. Gilbert also remembered that Grandin would not begin work or purchase the needed type from the foundry until `after Harris had promised to insure the payment for the printing.' Grandin's price to print five thousand copies was \$3,000, which would require Harris to impart essentially all of the property to which he had legal right. Printing began in September 1829.
			
			``JS likely dictated the text of this revelation sometime after the negotiations in June and before 25 August 1829, when Harris mortgaged his property to Grandin as payment for the publication, thus apparently fulfilling the revelation's injunction to `pay the printer's debt.' The language of the revelation suggests that Harris had already agreed to Grandin's terms but had not yet arranged payment. Grandin's brother-in-law later recalled that `Harris became for a time in some degree staggered in his confidence; but nothing could be done in the way of printing without his aid.' Once Harris mortgaged his property, however, Grandin considered himself paid in full. According to Gilbert, printing then proceeded: `As quick as Mr. Grandin got his type and got things all ready to commence the work, Hyrum Smith brought to the office 24 pages of manuscript on foolscap paper.'''
			
		% --
		\paragraph{HDDC}
		% --
		
			``It is unfortunate that a revelation of the magnitude of Section 19 received no introduction by the Prophet Joseph Smith in the \uline{History of the Church}. This revelation, which explains the doctrine of the atonement as it applies to the individual, was directed to Martin Harris...
			
			``Elder Orson Hyde...felt that it [this revelation] was a stinging rebuke to Martin Harris to be called a wicked man. Elder Hyde wrote an article about Martin harris in the \uline{Millennial Star} at a time when Martin had received a mission call from the James J. Strang church, of which he was then a member, to come to England as a missionary. Elder Hyde printed Section 19 with this article, and warned the Latter-day Saints in England of the unstable character of this man. In his scathing remarks, Elder Hyde suggests that since this revelation was the last Martin Harris was to receive (see \hyperref[DC19-32]{verse 32}), he would have to walk through the rest of his life without any inspiration from God.%
				\footnote{%
					% \label{}%
					HDDC footnote: ``\uline{Millennial Star} [Liverpool, England], November 15, 1846, pp. 124, 125.''
					}
			
			``The date of reception of this section is not known. In Table 19, all the entries that contain a date, agree it was given sometime in March 1830. It is interesting that this revelation follows \hyperref[DC20]{Section 20} in the \uline{History of the Church}.
			
			``There are three places within Section 19 in which major revisions have been made to the text as it is found in the Book of Commandments. These revisions are \hyperref[DC19-13]{verses 13--15}, \hyperref[DC19-20]{20}, and \hyperref[DC19-23]{23--25}. All other variations appear inconsequential'' (275--276).
			
	% -----
	\subsubsection{Versions}
	% -----
		\label{DC19-V}
	
		\begin{tabular}{p{1in}p{2in}p{1in}p{2in}}
			JSP		& \hyperref[EMS-1828-1829-JS-CIR-SUM-1829]{Circa Summer 1829}	& BC & Chapter 16, 39--42 \\
			DC 1835	& Section 44, 174--176											& TS & 3:24 (15 October 1842), 943--944 \\
			MS		& 4:8 (December 1843), 113--114									& DC 1844 & Section 44, 260--263 \\
		\end{tabular}
		
		\medskip
		
		See the \href{https://drive.google.com/file/d/1goAvNz8jS4R8slYfMG_db55Ty2z_vmgK/view?usp=sharing}{sections comparison} document for more information about the version history.

	% -----
	\subsubsection{Section 19:1} % *DC19-1 
	% -----
		\label{DC19-1}

			\footnote{%
				% \label{}%
				BC has for a heading, ``1 \textit{t:}.'' DC 1835 has, ``\textit{A commandment of God and not of man to Martin Harris, given (Manchester, New York, March. 1830,) by him who is eternal:}.'' DC 1844 has, ``\textit{A commandment of God and not of man to Martin Harris, given (Manchester, New York, March, 1830) by him who is eternal:}.''
				}%
		I AM Alpha and Omega,%
			\footnote{%
				% \label{}%
				See \hyperref[REVELATIONNT-1-8]{Revelation 1:8} and notes there.
				}
		Christ the Lord; yea, even I am he, the beginning and the end, the Redeemer of the world.

		\begin{framed}
		\scriptsize
			\begin{compactdesc}
				\item[JSP] YEA, even I, I am he, the beginning and the end: Yea, Alpha and Omega, Christ the Lord, the Redeemer of the world:
				\item[RB1] [\textit{missing}]
				\item[DC 1835] 1 I am Alpha and Omega, Christ the Lord; yea, even I am He, the beginning and the end, the Redeemer of the world: 
				\item[DC 1844] 1 I am Alpha and Omega, Christ the Lord; yea, even I am He, the beginning and the end, the Redeemer of the world: 
			\end{compactdesc}
		\end{framed}

	% -----
	\subsubsection{Section 19:2} % *DC19-2 
	% -----
		\label{DC19-2}

		I, having accomplished%
			\footnote{%
				% \label{}%
				I think this is the only place in the scriptures where ``accomplish'' occurs with ``will'' as an object. In the NT, ``accomplish'' appears most often in Luke, but it seems to mostly mean ``fulfill,'' though see \hyperref[LUKE-9-31]{Luke 9:31}.
				}
		and finished%
			\footnote{%
				% \label{}%
				Likewise with ``finish'' here, this is the only place where it occurs with ``will'' as an object, though see \hyperref[JOHN-4-34]{John 4:34}, which is close. The Greek word we should probably be looking to here is \hyperref[TELEIOO]{\grl{τελειόω}}.
				}
		the will of him whose I am, even the Father, concerning me---having done this that I might subdue all things unto myself---%
			\footnote{%
				% \label{}%
				See \hyperref[1CORINTHIANS-15-28]{1 Corinthians 15:28} and \hyperref[PHILIPPIANS-3-21]{Philippians 3:21}. There is a family of English words in the NT, such as ``subdue,'' ``submit,'' ``put under,'' and so on, which all translate the Greek \hyperref[HUPOTASSO]{\grl{ὑποτάσσω}}. See my note on ``submit'' at \hyperref[MOSIAH-3-19]{Mosiah 3:19}.
				}

		\begin{framed}
		\scriptsize
			\begin{compactdesc}
				\item[JSP] 2 I having accomplished and finished the will of him whose I am, even the Father: 3 Having done this, that I might subdue all things unto myself:
				\item[RB1] [\textit{missing}]
				\item[DC 1835] I having accomplished and finished the will of him whose I am, even the Father concerning me: having done this, that I might subdue all things unto myself: 
				\item[DC 1844] I having accomplished and finished the will of him whose I am, even the Father concerning me: having done this, that I might subdue all things unto myself: 
			\end{compactdesc}
		\end{framed}

	% -----
	\subsubsection{Section 19:3} % *DC19-3 
	% -----
		\label{DC19-3}

		Retaining all power,%
			\footnote{%
				% \label{}%
				Christ says that he has received ``all power'' in \hyperref[MATTHEW-28-18]{Matthew 28:18}.
				}
		even to the destroying of Satan and his works at the end of the world,
			\footnote{%
				% \label{}%
				What is the syntax here? I think ``the last great day of judgment'' is an object of the verb ``retaining,'' so it ought to read, ``Retaining the last great day of judgment.'' \hyperref[JOHN-5-22]{John 5:22} says that ``the Father judgeth no man, but hath committed all judgment unto the Son.''
				}%
		and the last great day of judgment, which I shall pass upon the inhabitants thereof, judging every man according to his works and the deeds%
			\footnote{%
				% \label{}%
				For ``works'' and ``deeds'' together, see \hyperref[PSALMS-28-4]{Psalms 28:4}, \hyperref[JEREMIAH-25-14]{Jeremiah 25:14}, \hyperref[JAMES-1-25]{James 1:25}, and \hyperref[HELAMAN-10-5]{Helaman 10:5}.
				}
		which he hath done.

		\begin{framed}
		\scriptsize
			\begin{compactdesc}
				\item[JSP] 4 Retaining all power, even to the destroying of satan and his works at the end of the world, and the last great day of judgment, which I shall pass upon the inhabitants thereof, judging every man according to his works, and the deeds which he hath done.
				\item[RB1] [\textit{missing}]
				\item[DC 1835] retaining all power, even to the destroying of satan and his works at the end of the world, and the last great day of judgment, which I shall pass upon the inhabitants thereof, judging every man according to his works, and the deeds which he hath done. 
				\item[DC 1844] retaining all power, even to the destroying of satan and his works at the end of the world, and the last great day of judgment, which I shall pass upon the inhabitants thereof, judging every man according to his works, and the deeds which he hath done. 
			\end{compactdesc}
		\end{framed}

	% -----
	\subsubsection{Section 19:4} % *DC19-4 
	% -----
		\label{DC19-4}

			\footnote{%
				% \label{}%
				This verse here begins an extremely important passage, ending with verse 12, about how to understand a certain concept in the scriptures. The principles here extend far beyond this passage. See my notes on the \hyperref[EPIST-REVELATION]{epistemology of revelation}, including more notes on these verses in the \hyperref[EPIST-REVELATION-DC]{DC notes}.
				}%
		And surely every man must repent or suffer, for%
			\footnote{%
				% \label{}%
				This ``for'' acts like a ``because'' here, and if you think of this verse in terms of an inference, it's really weird: ``I, God, am endless, and therefore every person must repent or suffer.'' Difficult to see the connection. You have to add in the next verse at least to get some sense out of it, but even then it's not that clear.
				}
		I, God, am endless.%
			\footnote{%
				\label{DC19-4-NOTE-endless}%
				For ``endless'' here and elsewhere, see \hyperref[MOSES-1-3]{Moses 1:3}, which is clearly influenced by this section. Moses 1 is around June 1830, and this section is around March 1830 (I'm not sure when I wrote this note, but this section is from the middle of 1829; there still could be an influence of this section on Moses, but they are not as close as I was assuming in time.)
				}

		\begin{framed}
		\scriptsize
			\begin{compactdesc}
				\item[JSP] 5 And surely every man must repent or suffer, for I God am endless:
				\item[RB1] [\textit{missing}]
				\item[DC 1835] And surely every man must repent or suffer, for I God am endless: 
				\item[DC 1844] And surely every man must repent or suffer, for I God am endless: 
			\end{compactdesc}
		\end{framed}

	% -----
	\subsubsection{Section 19:5} % *DC19-5 
	% -----
		\label{DC19-5}

		\hyperref[WHEREFORE]{Wherefore}, I revoke not the judgments which I shall pass, but woes shall go forth, weeping, wailing and gnashing of teeth, yea, to those who are found on my left hand.

		\begin{framed}
		\scriptsize
			\begin{compactdesc}
				\item[JSP] 6 Wherefore, I revoke not the judgments which I shall pass, but woes shall go forth, weeping, wailing and gnashing of teeth: 7 Yea, to those who are found on my left hand,
				\item[RB1] [\textit{missing}]
				\item[DC 1835] wherefore, I revoke not the judgments which I shall pass, but woes shall go forth, weeping, wailing and gnashing of teeth: yea, to those who are found on my left hand; 
				\item[DC 1844] wherefore, I revoke not the judgments which I shall pass, but woes shall go forth, weeping, wailing and gnashing of teeth: yea, to those who are found on my left hand: 
			\end{compactdesc}
		\end{framed}

	% -----
	\subsubsection{Section 19:6} % *DC19-6 
	% -----
		\label{DC19-6}

		Nevertheless, it is not written that there shall be no end to this torment,%
			\footnote{%
				% \label{}%
				Compare \hyperref[DC76-44]{DC 76:44--49}, especially verse 45, ``And the end thereof, neither the place thereof, nor their torment, no man knows,'' and 48, ``the end, the width, the height, the depth, and the misery thereof, they understand not, neither any man except those who are ordained unto this condemnation.'' 
				}
		but it is written \textit{endless torment}.%
			\footnote{%
				% \label{}%
				Besides its occurrence here, the phrase ``endless torment'' occurs seven other times in the scriptures, all of them in the BOM (see the list \hyperref[END-PHRASES]{here}). That in itself may be significant---the BOM is a modern revelation, and this could have been explained there, but it wasn't, for the reasons that the Lord discusses here.
				}

		\begin{framed}
		\scriptsize
			\begin{compactdesc}
				\item[JSP] nevertheless, it is not written, that there shall be no end to this torment; but it is written endless torment.
				\item[RB1] [\textit{missing}]
				\item[DC 1835] nevertheless it is not written, that there shall be no end to this torment; but it is written endless torment.
				\item[DC 1844] nevertheless it is not written, that there shall be no end to this torment; but it is written endless torment.
			\end{compactdesc}
		\end{framed}

	% -----
	\subsubsection{Section 19:7} % *DC19-7 
	% -----
		\label{DC19-7}

			\footnote{%
				% \label{}%
				For important commentary on this verse, see my notes on it at the \hyperref[EPIST-REVELATION-DC]{DC section} of the epistemology of revelation stuff.
				}%
		Again,%
			\footnote{%
				% \label{}%
				To me, this word here is the Lord saying, ``Here's another example of what I'm talking about.'' 
				}
		it is written \textit{eternal damnation};%
			\footnote{%
				% \label{}%
				This phrase, ``eternal damnation,'' occurs twice in the scriptures outside of this spot---\hyperref[MARK-3-29]{Mark 3:29} and \hyperref[DC29-44]{DC 29:44}. See my notes on the phrase \hyperref[ETERNITY-PHRASES]{here}.
				}
		wherefore it is more \hyperref[EXPRESS]{express}%
			\footnote{%
				% \label{}%
				For ``express,'' I think here we have the general meaning of ``plain, clear,'' but also meaning \#3 in the 1828, the pragmatic one, where the Lord is trying to accomplish a specific purpose---to get people to act---and so chooses these words for that purpose, even though they may not be exactly accurate.
				The adjectival use of ``express'' isn't common in the scriptures, but it does happen; see \hyperref[1SAMUEL-20-21]{1 Samuel 20:21} (adverb), \hyperref[EZEKIEL-1-3]{Ezekiel 1:3} (adverb), \hyperref[1TIMOTHY-4-1]{1 Timothy 4:1} (adverb), \hyperref[HEBREWS-1-3]{Hebrews 1:3}, \hyperref[MOSIAH-29-36]{Mosiah 29:36} (adverb), \hyperref[ALMA-9-24]{Alma 9:24} (adverb), \hyperref[DC107-43]{DC 107:43}, \hyperref[DC134-6]{134:6}, and \hyperref[DC138-40]{138:40}.
				}
		than other scriptures, that it might work upon the hearts of the children of men, altogether for my name's glory.

		\begin{framed}
		\scriptsize
			\begin{compactdesc}
				\item[JSP] 8 Again, it is written eternal damnation: wherefore it is more express than other scriptures, that it might work upon the hearts of the children of men, altogether for my name's glory:
				\item[RB1] [\textit{missing}]
				\item[DC 1835] 2 Again, it is written eternal damnation: wherefore it is more express than other scriptures, that it might work upon the hearts of the children of men, altogether for my name's glory: 
				\item[DC 1844] Again, it is written eternal damnation:---wherefore it is more express than other scriptures, that it might work upon the hearts of the children of men, altogether for my name's glory: 
			\end{compactdesc}
		\end{framed}

	% -----
	\subsubsection{Section 19:8} % *DC19-8 
	% -----
		\label{DC19-8}

		Wherefore, I will explain unto you this \hyperref[MYSTERY]{mystery}, for it is meet unto you to know even as mine apostles.%
		\footnote{%
			% \label{}%
			Woof---this is an occasion where you really wish that English nouns were declined, so we could see the case of ``apostles.'' The reason is that there are two options here: the ``apostles'' could refer to the ancient apostles, and so the verse would say, ``it is meet unto you to know even as my apostles in ancient times knew.'' But it could also be that ``apostles'' refers to JS and the other receivers of this revelation, so ``mine apostles'' would be like an explanation in the predicate of the identity of the people unto whom it's meet that they know. If we had declensions like Latin or Greek then we could tell which of these is right. The BC version, listed as the JSP here, doesn't really help much. The next verse, verse 9, might suggest that the ``apostles'' here are JS and the others.
			}

		\begin{framed}
		\scriptsize
			\begin{compactdesc}
				\item[JSP] 9 Wherefore, I will explain unto you, this mystery, for it is mete unto you, to know even as mine apostles.
				\item[RB1] [\textit{missing}]
				\item[DC 1835] wherefore, I will explain unto you, this mystery, for it is meet unto you, to know even as mine apostles. 
				\item[DC 1844] wherefore, I will explain unto you, this mystery, for it is meet unto you, to know even as mine apostles. 
			\end{compactdesc}
		\end{framed}

	% -----
	\subsubsection{Section 19:9} % *DC19-9 
	% -----
		\label{DC19-9}

		I speak unto you that are chosen in this thing, even as one, that you may enter into my rest.%
			\footnote{%
				% \label{}%
				See \hyperref[DC84-24]{DC 84:24}.
				}

		\begin{framed}
		\scriptsize
			\begin{compactdesc}
				\item[JSP] 10 I speak unto you that are chosen in this thing, even as one, that you may enter into my rest.
				\item[RB1] [\textit{missing}]
				\item[DC 1835] I speak unto you that are chosen in this thing, even as one, that you may enter into my rest. 
				\item[DC 1844] I speak unto you that are chosen in this thing, even as one, that you may enter into my rest. 
			\end{compactdesc}
		\end{framed}

	% -----
	\subsubsection{Section 19:10} % *DC19-10 
	% -----
		\label{DC19-10}

		For, behold, the \hyperref[MYSTERY]{mystery} of \hyperref[GODLINESS]{godliness},%
			\footnote{%
				\label{DC19-10-mystery}%
				See the \hyperref[GODLINESS-PHRASES]{occurrences} of this phrase, ``mystery of godliness,'' and my \hyperref[GODLINESS-mystery]{study} on it.
				}
		how great is it!
			\footnote{%
				% \label{}%
				Here begins the explanation of the phrases ``endless torment,'' ``eternal damnation,'' ``eternal punishment,'' and ``endless punishment.'' Their true meaning is that they refer to \textit{God's} punishment. They draw the ``endless''/``eternal'' adjectives from designations for God, or from God's attributes, according to this explanation.
				}%
		For, behold, I am endless, and the punishment which is given from my hand is endless punishment, for Endless is my name. \hyperref[WHEREFORE]{Wherefore}---

		\begin{framed}
		\scriptsize
			\begin{compactdesc}
				\item[JSP] 11 For behold, the mystery of Godliness how great is it? for behold I am endless, and the punishment which is given from my hand, is endless punishment, for endless is my name: 12 Wherefore---
				\item[RB1] [\textit{missing}]
				\item[DC 1835] For behold, the mystery of Godliness, how great is it? for behold I am endless, and the punishment which is given from my hand, is endless punishment, for endless is my name; wherefore---
				\item[DC 1844] For behold, the mystery of Godliness, how great is it? for behold I am endless, and the punishment which is given from my hand, is endless punishment, for endless is my name; wherefore---
			\end{compactdesc}
		\end{framed}

	% -----
	\subsubsection{Section 19:11} % *DC19-11 
	% -----
		\label{DC19-11}

		Eternal punishment is God's punishment.%
			\footnote{%
				% \label{}%
				Later, at \hyperref[DC19-20]{verse 20}, the Lord says that Martin Harris has experienced a very small taste of these punishments when the Lord withdrew the Spirit from him. That tells you something important about the nature of these punishments and what it would feel like to actually undergo them.
				}

		\begin{framed}
		\scriptsize
			\begin{compactdesc}
				\item[JSP] \phantom{.}
					
					\begin{tabular}{p{1in}>{\hangindent=1em}p{0.2in}p{2.0in}<{\raggedright}}
					Eternal punishment & \} & Endless punishment \\
					is God's punishment: & \} & is God's punishment:
					\end{tabular}
				
				\item[RB1] [\textit{missing}]
				\item[DC 1835] \phantom{.}
					
					\begin{tabular}{p{1in}>{\hangindent=1em}p{0.2in}p{2.0in}<{\raggedright}}
					Eternal punishment) & \phantom{.} & Endless punishment \\
					is God's punishment:) & \phantom{.} & is God's punishment:
					\end{tabular}
					
				\item[DC 1844] \phantom{.}
					
					\begin{tabular}{p{1in}>{\hangindent=1em}p{0.2in}p{2.0in}<{\raggedright}}
					Eternal punishment) & \phantom{.} & Endless punishment \\
					is God's punishment:) & \phantom{.} & is God's punishment:
					\end{tabular}
			\end{compactdesc}
		\end{framed}

	% -----
	\subsubsection{Section 19:12} % *DC19-12 
	% -----
		\label{DC19-12}

		Endless punishment is God's punishment.

		\begin{framed}
		\scriptsize
			\begin{compactdesc}
				\item[JSP] \phantom{.}
				
					\begin{tabular}{p{1in}>{\hangindent=1em}p{0.2in}p{2.0in}<{\raggedright}}
					Eternal punishment & \} & Endless punishment \\
					is God's punishment: & \} & is God's punishment:
					\end{tabular}
				
				\item[RB1] [\textit{missing}]
				\item[DC 1835] \phantom{.}
					
					\begin{tabular}{p{1in}>{\hangindent=1em}p{2.0in}p{2.0in}<{\raggedright}}
					Eternal punishment) & \phantom{.} & Endless punishment \\
					is God's punishment:) & \phantom{.} & is God's punishment:
					\end{tabular}
					
				\item[DC 1844] \phantom{.}
					
					\begin{tabular}{p{1in}>{\hangindent=1em}p{2.0in}p{2.0in}<{\raggedright}}
					Eternal punishment) & \phantom{.} & Endless punishment \\
					is God's punishment:) & \phantom{.} & is God's punishment:
					\end{tabular}
			\end{compactdesc}
		\end{framed}

	% -----
	\subsubsection{Section 19:13} % *DC19-13 
	% -----
		\label{DC19-13}

		Wherefore, I command you to repent, and keep the commandments which you have received by the hand of my servant Joseph Smith, Jun., in my name;

		\begin{framed}
		\scriptsize
			\begin{compactdesc}
				\item[JSP] ---
				\item[RB1] [\textit{missing}]
				\item[DC 1835] wherefore, I command you to repent, and keep the commandments which you have received by the hand of my servant Joseph Smith, jr. in my name: 
				\item[DC 1844] wherefore, I command you to repent, and keep the commandments which you have received by the hand of my servant Joseph Smith, jr. in my name: 
			\end{compactdesc}
		\end{framed}

	% -----
	\subsubsection{Section 19:14} % *DC19-14 
	% -----
		\label{DC19-14}

		And it is by my almighty power that you have received them;

		\begin{framed}
		\scriptsize
			\begin{compactdesc}
				\item[JSP] ---
				\item[RB1] [\textit{missing}]
				\item[DC 1835] and it is by my almighty power that you have received them: 
				\item[DC 1844] and it is by my almighty power that you have received them: 
			\end{compactdesc}
		\end{framed}

	% -----
	\subsubsection{Section 19:15} % *DC19-15 
	% -----
		\label{DC19-15}

		Therefore I command you to repent---repent, lest I smite you by the rod of my mouth,%
			\footnote{%
				% \label{}%
				This ``rod of my mouth'' bit comes from \hyperref[ISAIAH-11-4]{Isaiah 11:4}, and it also shows up at \hyperref[2NEPHI-21-4]{2 Nephi 21:4} and \hyperref[2NEPHI-30-9]{2 Nephi 30:9}. The Isaiah verse is talking about smiting the Earth, not an individual person.
				}
		and by my wrath, and by my anger, and your sufferings be sore---how sore you know not,%
			\footnote{%
				% \label{}%
				Notice here, in the JSP/BC version, the exclamation points, really underscoring what the Lord is saying here. It's very powerful. This passage, very unfortunately, does not appear in RB1; based on the patterns I've noticed there, though, my hypothesis would be that there were no exclamation points in that version, and that they were added when they went to print the BC. It is kind of interesting that we can't check though.
				}
		how exquisite%
			\footnote{%
				% \label{}%
				A rare word in the scriptures, appearing only four times---here, twice at \hyperref[ALMA-36-21]{Alma 36:21}, and once at \hyperref[JSHISTORY-1-31]{JS-H 1:31}. I'm willing to ignore the JS-H one because it would have been written so much later, and so the only real scriptural occurrences of ``exquisite'' are in the Alma verse. It reads: ``Yea, I say unto you, my son, that there can be nothing so exquisite and so bitter as was my pains. Yea, and again I say unto you, my son, that on the other hand there can be nothing so exquisite and sweet as was my joy.'' So Alma uses the word to describe both his pain \textit{and} his joy. The 1828 dictionary has these entries for ``exquisite'':
					\begin{compactitem}
						\item Literally, sought out or searched for with care; whence, choice; select. Hence,
						\item[1] Nice; exact; very excellent; complete; as a vase of \textit{exquisite} workmanship.
						\item[2] Nice; accurate; capable of nice perception; as \textit{exquisite} sensibility.
						\item[3] Nice; accurate; capable of nice discrimination; as \textit{exquisite} judgment, taste or discernment.
						\item[4] Being in the highest degree; extreme; as, to relish pleasure in an \textit{exquisite} degree. So we say, \textit{exquisite} pleasure or pain.
						\item[5] Very sensibly felt; as a painful and \textit{exquisite} impression on the nerves.
					\end{compactitem}
				The two meanings most relevant here seem to be \#4 and \#5. We do not understand the extreme nature of the suffering Christ endured during the events connected with the atonement.
				}
		you know not, yea, how hard to bear you know not.

		\begin{framed}
		\scriptsize
			\begin{compactdesc}
				\item[JSP] 13 Wherefore, I command you by my name, and by my Almighty power, that you repent: repent, lest I smite you by the rod of my mouth, and by my wrath, and by my anger, and your sufferings be sore:

					14 How sore you know not!

					15 How exquisite you know not!

					16 Yea, how hard to bear you know not!
				\item[RB1] [\textit{missing}]
				\item[DC 1835] therefore I command you to repent, repent, lest I smite you by the rod of my mouth, and by my wrath, and by my anger, and your sufferings be sore: how sore you know not! how exquisite you know not! yea, how hard to bear you know not! 
				\item[DC 1844] therefore I command you to repent, repent, lest I smite you by the rod of my mouth, and by my wrath, and by my anger, and your sufferings be sore: how sore you know not! how exquisit you know not! yea, how hard to bear you know not! 
			\end{compactdesc}
		\end{framed}

	% -----
	\subsubsection{Section 19:16} % *DC19-16 
	% -----
		\label{DC19-16}

		For behold, I, God, have suffered these things for all, that they might not suffer if they would repent;

		\begin{framed}
		\scriptsize
			\begin{compactdesc}
				\item[JSP] 17 For behold, I God have suffered these things for all, that they might not suffer, if they would repent,
				\item[RB1] [\textit{missing}]
				\item[DC 1835] For behold, I God have suffered these things for all, that they might not suffer, if they would repent, 
				\item[DC 1844] For behold, I God have suffered these things for all, that they might not suffer, if they would repent, 
			\end{compactdesc}
		\end{framed}

	% -----
	\subsubsection{Section 19:17} % *DC19-17 
	% -----
		\label{DC19-17}

		But if they would not repent they must suffer even as I;%
			\footnote{%
				% \label{}%
				``Even as I'' here could mean, ``They have to suffer, as I did,'' in the sense that they also have to suffer, but it could also mean ``They have to suffer, in the same way as I did.'' Note that Christ says he drank from the ``bitter cup'' in the next verse, and \hyperref[ALMA-40-26]{Alma 40:26} says that the wicked will drink the ``dregs of a bitter cup.'' The next verse also talks about Christ bleeding, and he says various times in the NT that his blood is ``of the new testament'' (\hyperref[MATTHEW-26-28]{Matthew 26:28}, \hyperref[MARK-14-24]{Mark 14:24}, \hyperref[LUKE-22-20]{Luke 22:20}). Think also of the food associated with the Passover, the ``bitter herbs'' (\hyperref[EXODUS-12-8]{Exodus 12:8} for example).
				}

		\begin{framed}
		\scriptsize
			\begin{compactdesc}
				\item[JSP] but if they would not repent, they must suffer even as I:
				\item[RB1] [\textit{missing}]
				\item[DC 1835] but if they would not repent, they must suffer even as I: 
				\item[DC 1844] but if they would not repent, they must suffer even as I: 
			\end{compactdesc}
		\end{framed}

	% -----
	\subsubsection{Section 19:18} % *DC19-18 
	% -----
		\label{DC19-18}

		Which suffering caused myself, even God, the greatest of all, to tremble because of pain, and to bleed%
			\footnote{%
				% \label{}%
				See \hyperref[LUKE-22-44]{Luke 22:44} and \hyperref[MOSIAH-3-7]{Mosiah 3:7}.
				}
		at every pore,%
			\footnote{%
				% \label{}%
				Notice the important difference between this version and the JSP/BC version: ``to bleed at every pore, both body and spirit,'' making it seem as though the bleeding was physical blood from the body as well as some kind of spiritual bleeding. If you think of spirit as a certain kind of matter then this is comprehensible and possibly even evidence for it.
				}
		and to suffer both body and spirit---and would%
			\footnote{%
				% \label{}%
				This ``would'' is the past/preterite form of ``will,'' and so it's saying, ``and I willed not to drink the bitter cup.'' In other words, Christ didn't wish to do it. 
				}
		that I might not drink the bitter cup,
			\footnote{%
				% \label{DC19-18-NOTE-shrink}%
				See \hyperref[1NEPHI-4-10]{1 Nephi 4:10} for the language here.
				}%
		and
			\footnote{%
				% \label{}%
				I think that ``shrink'' completes an implied use of ``might not,'' so it should read, ``I willed not to shrink.'' But there is also a reading that takes ``shrink'' as the result of drinking the bitter cup: ``I willed not to drink the bitter cup, with the result that I would shrink.'' 
				}%
		\hyperref[SHRINK]{shrink}---
			

		\begin{framed}
		\scriptsize
			\begin{compactdesc}
				\item[JSP] 18 Which suffering caused myself, even God, the greatest of all, to tremble because of pain, and to bleed at every pore, both body and spirit: 19 And would that I might not drink the bitter cup and shrink:
				\item[RB1] [\textit{missing}]
				\item[DC 1835] which suffering caused myself, even God, the greatest of all to tremble because of pain, and to bleed at every pore, and to suffer both body and spirit: and would that I might not drink the bitter cup and shrink: 
				\item[DC 1844] which suffering caused myself, even God, the greatest of all to tremble because of pain, and to bleed at every pore, and to suffer both body and spirit: and would that I might not drink the bitter cup and shrink: 
			\end{compactdesc}
		\end{framed}

	% -----
	\subsubsection{Section 19:19} % *DC19-19 
	% -----
		\label{DC19-19}

		Nevertheless, glory be to the Father, and I partook and finished%
			\footnote{%
				% \label{}%
				See \hyperref[DC76-69]{DC 76:69} and my notes there.
				}
		my preparations%
			\footnote{%
				% \label{}%
				I connected this word ``preparations'' with the ``preparation'' for the passover, spoken of in various places (\hyperref[JOHN-19-14]{John 19:14} for example). The Greek word is \gr{παρασκευή}. For that word, see BDAG 685a: ``primary sense `preparation'...in our literature only of a definite day, as the \textit{day of preparation} for a festival; according to Israel's usage...it was Friday, on which day everything had to be prepared for the Sabbath, when no work was permitted (\hyperref[MATTHEW-27-62]{Matthew 27:62}).'' Whatever Christ is talking about here---the prayer in the garden or something else---in this revelation he describes that as a series of ``preparations,'' connecting to the ``preparations'' that people made for the special festivals they celebrated, and most important of all the Passover.
				}
		unto the children of men.

		\begin{framed}
		\scriptsize
			\begin{compactdesc}
				\item[JSP] 20 Nevertheless, glory be to the Father, and I partook and finished my preparations unto the children of men:
				\item[RB1] [\textit{missing}]
				\item[DC 1835] nevertheless, glory be to the Father, and I partook and finished my preparations unto the children of men: 
				\item[DC 1844] nevetheless, glory be to the Father, and I partook and finished my preparations unto the children of men: 
			\end{compactdesc}
		\end{framed}

	% -----
	\subsubsection{Section 19:20} % *DC19-20 
	% -----
		\label{DC19-20}

		Wherefore, I command you again to repent, lest I humble you with my almighty power; and that you confess your sins, lest you suffer these punishments of which I have spoken, of which in the smallest, yea, even in the least degree you have tasted at the time I withdrew my Spirit.

		\begin{framed}
		\scriptsize
			\begin{compactdesc}
				\item[JSP] 21 Wherefore, I command you again by my Almighty power, that you confess your sins, lest you suffer these punishments of which I have spoken, of which in the smallest, yea, even in the least degree you have tasted at the time I withdrew my Spirit.
				\item[RB1] By my Almighty power that ye confess your sins lest ye suffer these punishments of which I have spoken of which in the smalest yea in the least decree ye have tasted at a time I withdrew my spirit
				\item[DC 1835] wherefore, I command you again to repent lest I humble you by my almighty power, and that you confess your sins lest you suffer these punishments of which I have spoken, of which in the smallest, yea, even in the least degree you have tasted at the time I withdrew my Spirit. 
				\item[DC 1844] wherefore, I command you again to repent lest I humble you with my almighty power, and that you confess your sins lest you suffer these punishments of which I have spoken, of which in the smallest, yea, even in the least degree you have tasted at the time I withdrew my Spirit. 
			\end{compactdesc}
		\end{framed}

	% -----
	\subsubsection{Section 19:21} % *DC19-21 
	% -----
		\label{DC19-21}

		And I command you that you preach naught but repentance, and show not these things unto the world until it is wisdom in me.

		\begin{framed}
		\scriptsize
			\begin{compactdesc}
				\item[JSP] 22 And I command you, that you preach nought but repentance; and show not these things, neither speak these things unto the world,
				\item[RB1] I command you that ye preach naught but repentance \& shew not these things neither speak these things unto the World.
				\item[DC 1835] And I command you, that you preach nought but repentance; and show not these things unto the world until it is wisdom in me; 
				\item[DC 1844] And I command you, that you preach nought but repentance; and show not these things unto the world until it is wisdom in me; 
			\end{compactdesc}
		\end{framed}

	% -----
	\subsubsection{Section 19:22} % *DC19-22 
	% -----
		\label{DC19-22}

		For they cannot bear meat now, but milk they must receive;%
			\footnote{%
				% \label{}%
				See \hyperref[1CORINTHIANS-3-2]{1 Corinthians 3:2} and \hyperref[HEBREWS-5-12]{Hebrews 5:12--14}, but also \hyperref[DC109-14]{DC 109:14--15}.
				}
		wherefore, they must not know these things, \hyperref[LEST]{lest} they perish.%
			\footnote{%
				% \label{}%
				Remember that ``lest'' means ``so that not.'' For some reason my mind wants to read the end of this verse as: ``they can't know these things, because what if they die?'' But that is not what it means, I think. It should be---``they must not know these things, so that they do not perish.'' There's something about \textit{knowing} the things that plays a role in the people \textit{perishing}. Could be referring to how, by acquiring greater knowledge, they become more accountable, and so are more susceptible to losing the Spirit when they don't live by that knowledge.
				}

		\begin{framed}
		\scriptsize
			\begin{compactdesc}
				\item[JSP] for they can not bear meat, but milk they must receive: 23 Wherefore, they must not know these things lest they perish:
				\item[RB1] for they cannot bear meat but milk they must receive. Wherefore, they must not know these things lest they perish
				\item[DC 1835] for they cannot bear meat now, but milk they must receive: wherefore, they must not know these things lest they perish: 
				\item[DC 1844] for they cannot bear meat now, but milk they must receive: wherefore, they must not know these things lest they perish: 
			\end{compactdesc}
		\end{framed}

	% -----
	\subsubsection{Section 19:23} % *DC19-23 
	% -----
		\label{DC19-23}

		Learn of me, and listen to my words; walk in the \hyperref[MEEKNESS]{meekness} of my Spirit, and you shall have peace%
			\footnote{%
				% \label{}%
				See \hyperref[JOHN-14-27]{John 14:27}.
				}
		in me.%
			\footnote{%
				% \label{}%
				Both the JSP/BC and RB1 are different here---they connect verses 23 and 24 immediately together as part of the same thought, so the current version has a significant change here.
				}

		\begin{framed}
		\scriptsize
			\begin{compactdesc}
				\item[JSP] 24 Wherefore, learn of me, and listen to my words; walk in the meekness of my Spirit and you shall have peace in me,
				\item[RB1] Wherefore learn of me \& listen to my words walk in the meekness of my spirit \& ye shall have peace in me
				\item[DC 1835] learn of me, and listen to my words; walk in the meekness of my Spirit and you shall have peace in me: 
				\item[DC 1844] learn of me, and listen to my words; walk in the meekness of my Spirit and you shall have peace in me: 
			\end{compactdesc}
		\end{framed}

	% -----
	\subsubsection{Section 19:24} % *DC19-24 
	% -----
		\label{DC19-24}

		I am Jesus Christ; I came by the will of the Father, and I do his will.

		\begin{framed}
		\scriptsize
			\begin{compactdesc}
				\item[JSP] Jesus Christ by the will of the Father.
				\item[RB1] (Jesus Christ) by the will of the Father
				\item[DC 1835] I am Jesus Christ: I came by the will of the Father, and I do his will.
				\item[DC 1844] I am Jesus Christ; I came by the will of the Father, and I do his will.
			\end{compactdesc}
		\end{framed}

	% -----
	\subsubsection{Section 19:25} % *DC19-25 
	% -----
		\label{DC19-25}

		And again, I command thee that thou shalt not covet thy neighbor's wife; nor seek thy neighbor's life.%
			\footnote{%
				% \label{}%
				Could be that these were things Martin in particular had a problem with; he might have been into his neighbor's wife, and maybe he had thought about harming that neighbor, or someone else?
				}

		\begin{framed}
		\scriptsize
			\begin{compactdesc}
				\item[JSP] 25 And again: I command you, that thou shalt not covet thy neighbor's wife. 26 Nor seek thy neighbor's life.
				\item[RB1] \& again I command you that thou shalt not covet thy Neibours wife. nor seek thy Neibours life.
				\item[DC 1835] 3 And again: I command thee, that thou shalt not covet thy neighbor's wife. Nor seek thy neighbor's life. 
				\item[DC 1844] 3 And again, I command thee, that thou shalt not covet thy neighbor's wife. Nor seek thy neighbor's life. 
			\end{compactdesc}
		\end{framed}

	% -----
	\subsubsection{Section 19:26} % *DC19-26 
	% -----
		\label{DC19-26}

		And again, I command thee that thou shalt not covet thine own property, but impart it freely to the printing of the Book of Mormon, which contains the truth%
			\footnote{%
				% \label{}%
				The RB1 version lacks ``the truth.''
				}
		and the word of God---

		\begin{framed}
		\scriptsize
			\begin{compactdesc}
				\item[JSP] 27 And again: I command you, that thou shalt not covet thine own property, but impart it freely to the printing of the book of Mormon, which contains the truth and the word of God,
				\item[RB1] \& again I command you that thou shalt not covet thine own Property but impart it freely to the printing of the Book\sout{s} of Mormon which contains of the word of God
				\item[DC 1835] And again: I command thee, that thou shalt not covet thine own property, but impart it freely to the printing of the book of Mormon, which contains the truth and the word of God, 
				\item[DC 1844] And again: I command thee, that thou shalt not covet thine own property, but impart it freely to the printing of the book of Mormon, which contains the truth and the word of God, 
			\end{compactdesc}
		\end{framed}

	% -----
	\subsubsection{Section 19:27} % *DC19-27 
	% -----
		\label{DC19-27}

		Which is my word to the Gentile, that soon it may go to the Jew, of whom the Lamanites are a remnant, that they may believe the gospel, and look not for a Messiah to come who has already come.

		\begin{framed}
		\scriptsize
			\begin{compactdesc}
				\item[JSP] which is my word to Gentile, that soon it may go to the Jew, of which the Lamanites are a remnant; that they may believe the gospel, and look not for a Messiah to come which has already come.
				\item[RB1] which is my word to Gentile that soon it may go to the Jew of which the Lamanites \sout{of} are a Remnant that they may Believe the Gospel \& look not for a Masiah to come which has already come
				\item[DC 1835] which is my word to the Gentile, that soon it may go to the Jew, of whom the Lamanites are a remnant: that they may believe the gospel, and look not for a Messiah to come who has already come.
				\item[DC 1844] which is my word to the Gentile, that soon it may go to the Jew, of whom the Lamanites are a remnant; that they may believe the gospel, and look not for a Messiah to come who has already come.
			\end{compactdesc}
		\end{framed}

	% -----
	\subsubsection{Section 19:28} % *DC19-28 
	% -----
		\label{DC19-28}

		And again, I command thee that thou shalt pray vocally as well as in thy heart;%
			\footnote{%
				% \label{}%
				Here both JSP/BC and RB1 have the instruction to pray ``to thyself,'' instead of ``in thy heart.''
				}
		yea, before the world as well as in secret, in public as well as in private.

		\begin{framed}
		\scriptsize
			\begin{compactdesc}
				\item[JSP] 28 And again: I command you, that thou shalt pray vocally as well as to thyself: 29 Yea, before the world as well as in secret; in public as well as in private.
				\item[RB1] \& again I command you that thou shalt pray vocally as well as to thyself yea before the World as well as in seecret in Publick as well as in private
				\item[DC 1835] 4 And again. I command thee, that thou shalt pray vocally as well as in thy heart: yea, before the world as well as in secret; in public as well as in private. 
				\item[DC 1844] 4 And again, I command thee, that thou shalt pray vocally as well as in thy heart: yea, before the world as well as in secret; in public as well as in private. 
			\end{compactdesc}
		\end{framed}

	% -----
	\subsubsection{Section 19:29} % *DC19-29 
	% -----
		\label{DC19-29}

		And thou shalt declare glad tidings, yea, publish it upon the mountains, and upon every high place, and among every people that thou shalt be permitted to see.

		\begin{framed}
		\scriptsize
			\begin{compactdesc}
				\item[JSP] 30 And thou shalt declare glad tidings; yea, publish it upon the mountains, and upon every high place, and among every people which thou shalt be permitted to see.
				\item[RB1] \& thou shalt declare glad tidings yea Publish it upon the Mountains \& upon evry high place \& among evry People which thou shalt be permited to see
				\item[DC 1835] And thou shalt declare glad tidings; yea, publish it upon the mountains, and upon every high place, and among every people that thou shalt be permitted to see. 
				\item[DC 1844] And thou shalt declare glad tidings: yea, publish it to the mountains, and upon every high place, and among every people that thou shalt be permitted to see. 
			\end{compactdesc}
		\end{framed}

	% -----
	\subsubsection{Section 19:30} % *DC19-30 
	% -----
		\label{DC19-30}

		And thou shalt do it with all humility, trusting in me, reviling not against revilers.

		\begin{framed}
		\scriptsize
			\begin{compactdesc}
				\item[JSP] 31 And thou shalt do it with all humility, trusting in me, reviling not against revilers.
				\item[RB1] \& thou shalt do it with all humility trusting in me Reviling not against Revilers
				\item[DC 1835] And thou shalt do it with all humility, trusting in me, reviling not against revilers. 
				\item[DC 1844] And thou shalt do it with all humility, trusting in me, reviling not against revilers.---
			\end{compactdesc}
		\end{framed}

	% -----
	\subsubsection{Section 19:31} % *DC19-31 
	% -----
		\label{DC19-31}

		And of tenets%
		\footnote{%
			% \label{}%
			This is the only use of ``tenet'' or ``tenets'' in the scriptures, though it does appear once at \hyperref[JSHISTORY-1-9]{JS-H 1:9}. The 1828 dictionary has this: ``Any opinion, principle, dogma or doctrine which a person believes or maintains as true; as the tenets of Plato or of Cicero. The tenets of Christians are adopted from the Scriptures; but different interpretations give rise to a great diversity of tenets.''
			} 
		thou shalt not talk, but thou shalt declare repentance and faith on the Savior, and remission of sins by baptism, and by fire, yea, even the Holy Ghost.%
			\footnote{%
				% \label{}%
				This list of items seems to foreshadow what JS later writes in the \hyperref[ARTICLESOFFAITH]{fourth article of faith}.
				} 

		\begin{framed}
		\scriptsize
			\begin{compactdesc}
				\item[JSP] 32 And of tenets thou shalt not talk, but thou shalt declare repentance and faith on the Savior and remission of sins by baptism and by fire; yea, even the Holy Ghost.
				\item[RB1] \& of tenets thou shalt not \sout{talk} talk but thou shalt declare Repentance \& faith on the Saveiour \& Remissions of Sins by Baptism \& by fire yea even the Holy Ghost
				\item[DC 1835] And of tenets thou shalt not talk, but thou shalt declare repentance and faith on the Savior, and remission of sins by baptism and by fire; yea, even the Holy Ghost.
				\item[DC 1844] And of tenets thou shalt not talk, but thou shalt declare repentance and faith on the Savior, and remission of sins by baptism and by fire; yea, even the Holy Ghost.
			\end{compactdesc}
		\end{framed}

	% -----
	\subsubsection{Section 19:32} % *DC19-32 
	% -----
		\label{DC19-32}

		Behold, this is a great and the last commandment which I shall give unto you concerning this matter; for this shall suffice for thy daily walk, even unto the end of thy life.

		\begin{framed}
		\scriptsize
			\begin{compactdesc}
				\item[JSP] 33 Behold this is a great and the last commandment which I shall give unto you: 34 For this shall suffice for thy daily walk even unto the end of thy life.
				\item[RB1] Behold this is a great \& last commandment which I shall give unto you for this shall suffice for thy daily walk even unto the end of thy life
				\item[DC 1835] 5 Behold, this is a great, and the last commandment which I shall give unto you concerning this matter: for this shall suffice for thy daily walk even unto the end of thy life. 
				\item[DC 1844] 5 Behold this is a great, and the last commandment which I shall give unto you concerning this matter: for this shall suffice for thy daily walk even unto the end of thy life. 
			\end{compactdesc}
		\end{framed}

	% -----
	\subsubsection{Section 19:33} % *DC19-33 
	% -----
		\label{DC19-33}

		And misery thou shalt receive if thou wilt slight these counsels, yea, even the destruction of thyself and property.

		\begin{framed}
		\scriptsize
			\begin{compactdesc}
				\item[JSP] 35 And misery thou shalt receive, if thou wilt slight these counsels; Yea, even destruction of thyself and property.
				\item[RB1] \& misery shalt thou Receive if thou wilt sleight these Councils yea even distruction of thyself \& Property.
				\item[DC 1835] And misery thou shalt receive, if thou wilt slight these counsels; yea, even destruction of thyself and property. 
				\item[DC 1844] And misery thou shalt receive, if thou wilt slight these counsels; yea, even the destruction of thyself and property. 
			\end{compactdesc}
		\end{framed}

	% -----
	\subsubsection{Section 19:34} % *DC19-34 
	% -----
		\label{DC19-34}

		Impart a portion of thy property, yea, even part of thy lands, and all save the support of thy family.

		\begin{framed}
		\scriptsize
			\begin{compactdesc}
				\item[JSP] 36 Impart a portion of thy property; Yea, even a part of thy lands and all save the support of thy family.
				\item[RB1] Impart a portion of thy Property yea even a part of thy lands \& all save the support of thy family
				\item[DC 1835] Impart a portion of thy property; yea, even part of thy lands and all save the support of thy family. 
				\item[DC 1844] Impart a portion of thy property; yea, even part of thy lands and all save the support of thy family. 
			\end{compactdesc}
		\end{framed}

	% -----
	\subsubsection{Section 19:35} % *DC19-35 
	% -----
		\label{DC19-35}

		Pay the debt thou hast contracted with the printer. Release thyself from bondage.

		\begin{framed}
		\scriptsize
			\begin{compactdesc}
				\item[JSP] 37 Pay the printer's debt. 38 Release thyself from bondage.
				\item[RB1] Pay the Printers debt Release thyself from Bondage
				\item[DC 1835] Pay the debt thou hast contracted with the printer. Release thyself from bondage. 
				\item[DC 1844] Pay the debt thou hast contracted with the printer. Release thyself from bondage. 
			\end{compactdesc}
		\end{framed}

	% -----
	\subsubsection{Section 19:36} % *DC19-36 
	% -----
		\label{DC19-36}

		Leave thy house and home, except when thou shalt desire to see thy family;

		\begin{framed}
		\scriptsize
			\begin{compactdesc}
				\item[JSP] 39 Leave thy house and home, except when thou shalt desire to see them.
				\item[RB1] Leave thy House \& home except when thou shalt desire to see them
				\item[DC 1835] Leave thy house and home, except when thou shalt desire to see thy family. 
				\item[DC 1844] Leave thy house and home, except when thou shalt desire to see thy family. 
			\end{compactdesc}
		\end{framed}

	% -----
	\subsubsection{Section 19:37} % *DC19-37 
	% -----
		\label{DC19-37}

		And speak freely to all; yea, preach, exhort, declare the truth, even with a loud voice, with a sound of rejoicing, crying---Hosanna, hosanna, blessed be the name of the Lord God!

		\begin{framed}
		\scriptsize
			\begin{compactdesc}
				\item[JSP] 40 And speak freely to all: Yea, preach, exhort, declare the truth, even with a loud voice; with a sound of rejoicing, crying hosanna! hosanna! blessed be the name of the Lord God.
				\item[RB1] \& speak freely to all yea Preach exhort declare the truth even With a loud voice with a sound of rejoiceing crying Hozannah Hozannah blesed be the name of the Lord God
				\item[DC 1835] And speak freely to all: yea, preach, exhort, declare the truth, even with a loud voice; with a sound of rejoicing, crying hosanna! hosanna! blessed be the name of the Lord God.
				\item[DC 1844] And speak freely to all: yea, preach, exhort, declare the truth, even with a loud voice; with a sound of rejoicing, crying hosanna! hosanna! blessed be the name of the Lord God.
			\end{compactdesc}
		\end{framed}

	% -----
	\subsubsection{Section 19:38} % *DC19-38 
	% -----
		\label{DC19-38}

		Pray always, and I will pour out my Spirit upon you, and great shall be your blessing---yea, even more than if you should obtain treasures of earth and corruptibleness to the extent thereof.

		\begin{framed}
		\scriptsize
			\begin{compactdesc}
				\item[JSP] 41 Pray always and I will pour out my Spirit upon you, and great shall be your blessing: 42 Yea, even more than if you should obtain treasures of earth, and corruptibleness to the extent thereof.
				\item[RB1] pray always \& I will pour out my spirit upon you \& great shall be your Blessing yea even more then if you should Obtain treasures of Earth \& corruptableness to the extent thereof
				\item[DC 1835] 6 Pray always and I will pour out my Spirit upon you, and great shall be your blessing: yea, even more than if you should obtain treasures of earth, and corruptibleness to the extent thereof. 
				\item[DC 1844] 6 Pray always and I will pour out my Spirit upon you, and great shall be your blessing: yea, even more than if you should obtain treasures of earth and corruptibleness to the extent thereof. 
			\end{compactdesc}
		\end{framed}

	% -----
	\subsubsection{Section 19:39} % *DC19-39 
	% -----
		\label{DC19-39}

		Behold, canst thou read this without rejoicing and lifting up thy heart for gladness?

		\begin{framed}
		\scriptsize
			\begin{compactdesc}
				\item[JSP] 43 Behold, canst thou read this without rejoicing, and lifting up thy heart for gladness;
				\item[RB1] Behold canst thou read this without Rejoiceing \& lifting up thy heart for Gladness
				\item[DC 1835] Behold, canst thou read this without rejoicing and lifting up thy heart for gladness; 
				\item[DC 1844] Behold, canst thou read this without rejoicing and lifting up thy heart for gladness; 
			\end{compactdesc}
		\end{framed}

	% -----
	\subsubsection{Section 19:40} % *DC19-40 
	% -----
		\label{DC19-40}

		Or canst thou run about longer as a blind guide?%
			\footnote{%
				% \label{}%
				So Harris is the blind guide? Guiding who?
				}

		\begin{framed}
		\scriptsize
			\begin{compactdesc}
				\item[JSP] or canst thou run about longer as a blind guide;
				\item[RB1] or canst thou run about longer as a blind guide
				\item[DC 1835] or canst thou run about longer as a blind guide; 
				\item[DC 1844] or canst thou run about longer as a blind guide; 
			\end{compactdesc}
		\end{framed}

	% -----
	\subsubsection{Section 19:41} % *DC19-41 
	% -----
		\label{DC19-41}

		Or canst thou be humble and meek, and conduct thyself wisely before me?%
			\footnote{%
				% \label{}%
				So there are two options: be blind, or be humble and meek and conduct yourself wisely---if you do all of that, you won't be blind.
				}
		Yea, come unto me thy Savior. Amen.

		\begin{framed}
		\scriptsize
			\begin{compactdesc}
				\item[JSP] or canst thou be humble and meek and conduct thyself wisely before me: 44 Yea, come unto me thy Savior. Amen.
				\item[RB1] or canst thou <be> humble \& meek \& conduct thyself wisely before me yea come unto me thy Saveiour amen
				\item[DC 1835] or canst thou be humble and meek and conduct thyself wisely before me: yea, come unto me thy Savior. Amen.
				\item[DC 1844] or canst thou be humble and meek and conduct thyself wisely before me: yea, come unto me thy Savior: Amen.
			\end{compactdesc}
		\end{framed}